% !TEX root = ../main.tex
\chapter{型・値・関数}
\label{chap:types-functions}

\begin{chapterabstract}
本章では、Lean 4を「証明を書くためだけの特殊な道具」としてではなく、まず型と関数を扱う小さな言語として導入する。
ソフトウェアエンジニアにとって、型、値、関数は日常的な概念である。
しかしLeanでは、それらが後で仕様や証明を組み立てるための材料にもなる。
本章では、\texttt{Nat}、\texttt{Bool}、\texttt{String}、\texttt{List}、\texttt{Option}を使い、\texttt{def}による関数定義、\texttt{\#eval}による実行確認、関数合成による小さなパイプラインを扱う。
圏論の用語はまだ前面に出さないが、「型を対象、関数を変換、合成を処理の連結として見る」ための準備を行う。
\end{chapterabstract}

\begin{learninggoals}
\item Lean 4における型と値を、実務上の「値に対する約束」として読める。
\item \texttt{Nat}、\texttt{Bool}、\texttt{String}、\texttt{List}、\texttt{Option}の基本的な使い方を理解できる。
\item \texttt{def}で小さな関数を定義し、\texttt{\#eval}で動作を確認できる。
\item 関数をつないで、処理パイプラインとして読む感覚を持てる。
\item 後続章で扱う仕様、証明、圏論の語彙へ進むための最小限のLeanコードを読める。
\end{learninggoals}

\section{この章を読む理由}

ソフトウェア開発では、値の形が変わるたびに小さな約束が生まれる。
ユーザーIDは数値なのか文字列なのか。
検索結果は必ず存在するのか、それとも存在しないことがあるのか。
APIから返る値は単体なのか、リストなのか。
このような約束が曖昧なままだと、実装は動いても、仕様の読み方は不安定になる。

Leanで最初に確認するのは、まさにこの約束である。
Leanは、値に必ず型を与える。
そして関数には、入力の型と出力の型を与える。
そのため、Leanのコードを読むときは、「この処理は何を受け取り、何を返すのか」を常に意識することになる。

\begin{intuitionbox}
型は、値の置き場所というよりも、値についての約束である。
たとえば\texttt{Nat}という型は「この値は自然数として扱う」という約束であり、\texttt{Option String}という型は「文字列があるかもしれないし、ないかもしれない」という約束である。
この約束を明示することが、後で仕様を命題として書くための第一歩になる。
\end{intuitionbox}

\FigurePlaceholder{fig-ch01-types-values-functions.png}{0.90}{型・値・関数の関係}{fig:ch01-types-values-functions}

\section{ソフトウェア上の問題：値の約束が曖昧になる}

実務では、関数の名前やコメントから意図を読み取ることが多い。
しかし、名前やコメントだけでは機械的には検査されない。
たとえば、次のような状況を考える。

\begin{itemize}
\item ユーザー名が空文字列かもしれない。
\item ユーザー検索は失敗するかもしれない。
\item 正規化済みの文字列と未正規化の文字列が同じ型で混ざっている。
\item 単一の値を返す処理と、リストを返す処理が混同される。
\end{itemize}

このような曖昧さは、テストで一部は発見できる。
しかし、関数の境界で「何が入力として許されるか」「何が出力として返るか」を型に書いておくと、そもそも混ぜてはいけない値を混ぜにくくなる。
Leanでは、この型の情報が非常に中心的な役割を持つ。

\begin{softwarebox}
本章で行うことは、本番アプリケーションをLeanで書き直すことではない。
実務コードの中から、仕様の核になりそうな小さな変換を取り出し、型と関数として書いてみることである。
たとえば「入力文字列を正規化する」「IDからユーザー名を探す」「見つからない場合を\texttt{Option}で表す」といった小さな処理が対象になる。
\end{softwarebox}


\section{型注釈と\texttt{\#check}}

Leanでは、すべての式に型がある。
数値リテラル、文字列リテラル、真偽値、リスト、関数、型そのものにも型がある。
最初は、「式を見たら、その式の型を確認する」という読み方に慣れることが大切である。

型を明示したいときは、次の形で書く。

\begin{center}
\texttt{(式 : 型)}
\end{center}

この\texttt{:}は、「左の式を右の型として読む」という注釈である。
日本語で読むなら、\texttt{(3 : Nat)}は「\texttt{3}を\texttt{Nat}として扱う」と読める。
ここでの\texttt{:}は、代入ではない。
値を変換しているのでもない。
Leanに対して、その式の型を明示しているだけである。

\begin{definitionbox}
本章では、型、値、式を次のように区別する。

\begin{itemize}
\item 型は、式がどの種類の値として扱われるかを表す。
\item 値は、計算の結果として得られる具体的なものを表す。
\item 式は、値そのもの、または値を作るための書き方を表す。
\end{itemize}

たとえば\texttt{3 + 4}は式であり、評価すると\texttt{7}という値になる。
その式は\texttt{Nat}型として扱える。
\end{definitionbox}

型を確認するには、\texttt{\#check}を使う。
\texttt{\#check}は、Leanに「この名前や式の型は何か」と尋ねるコマンドである。

\begin{listing}[htbp]
\begin{minted}{lean4}
#check Nat
#check Bool
#check String

#check (3 : Nat)
#check (true : Bool)
#check ("lean" : String)

#check List Nat
#check Option String
\end{minted}
\caption{型注釈と型確認}
\label{lst:ch01-type-annotation-check}
\end{listing}

このコードでは、\texttt{Nat}、\texttt{Bool}、\texttt{String}という型名も\texttt{\#check}している。
Leanでは、型名もまた式として扱われる。
たとえば\texttt{Nat}は、自然数の値そのものではなく、自然数の型を表す名前である。

\begin{leanbox}
\texttt{\#check Nat}を実行すると、Leanは\texttt{Nat : Type}のような結果を返す。
これは「\texttt{Nat}は\texttt{Type}に属する」と読める。
\texttt{Type}は、型を集めた階層の入口である。
宇宙レベルなどの詳しい話は、本章では扱わない。
今は「型にも型がある」とだけ読めば十分である。
\end{leanbox}

\texttt{List Nat}や\texttt{Option String}は、二つの単語が並んでいるように見える。
これは、\texttt{List}や\texttt{Option}に型を渡して、新しい型を作っている。
Leanでは、型に型を渡すときにも、通常の関数適用と同じように空白で並べる。
\texttt{List Nat}は「\texttt{Nat}のリスト型」、\texttt{Option String}は「\texttt{String}があるかもしれない型」と読める。

\begin{warningbox}
\texttt{(3 : Nat)}のような型注釈は、Leanに型を教えるための書き方である。
合わない型を無理に付けることはできない。
たとえば、文字列\texttt{"lean"}を\texttt{Nat}として扱うことはできない。
型注釈は変換ではなく、確認と明示のための構文である。
\end{warningbox}

\section{基本型と値の書き方}

この節では、Leanの基本的な型と、その値の書き方を確認する。
ここで扱うのは、\texttt{Nat}、\texttt{Bool}、\texttt{String}、\texttt{List}、\texttt{Option}である。
いずれも、以後の章で繰り返し使う。

\begin{center}
\begin{tabular}{lll}
\toprule
型 & 値の例 & 読み方 \\
\midrule
\texttt{Nat} & \texttt{0}, \texttt{1}, \texttt{2 + 3} & 自然数 \\
\texttt{Bool} & \texttt{true}, \texttt{false} & 真偽値 \\
\texttt{String} & \texttt{"lean"} & 文字列 \\
\texttt{List Nat} & \texttt{[1, 2, 3]} & 自然数のリスト \\
\texttt{Option Nat} & \texttt{some 3}, \texttt{none} & 自然数がある、またはない \\
\bottomrule
\end{tabular}
\end{center}

\texttt{Nat}の値は、\texttt{0}、\texttt{1}、\texttt{2}のように書く。
足し算は\texttt{+}で書ける。
本章では自然数の基本的な計算だけを使い、整数や小数には進まない。

\texttt{Bool}の値は、\texttt{true}と\texttt{false}である。
\texttt{\&\&}は論理積を表す。
\texttt{true \&\& false}は、両方が真のときだけ真になる式である。

\texttt{String}の値は、二重引用符で囲んで書く。
文字列の連結には\texttt{++}を使える。
\texttt{"lean" ++ "4"}は、二つの文字列をつないだ文字列になる。

\texttt{List A}は、\texttt{A}型の値を並べた型である。
角括弧\texttt{[...]}を使ってリストを書く。
空リスト\texttt{[]}だけでは要素型が分からないことがあるため、必要に応じて\texttt{([] : List Nat)}のように型注釈を付ける。

\texttt{Option A}は、\texttt{A}型の値がある場合と、ない場合を表す型である。
値がある場合は\texttt{some a}、値がない場合は\texttt{none}と書く。
\texttt{none}だけでは何の\texttt{Option}か分からないことがあるため、\texttt{(none : Option Nat)}のように型注釈を付ける。

\begin{listing}[htbp]
\begin{minted}{lean4}
#eval (3 : Nat) + 4
#eval true && false
#eval "lean" ++ "4"

#eval ([1, 2, 3] : List Nat).length
#eval ([true, false] : List Bool)

#eval (some 3 : Option Nat)
#eval (none : Option Nat)
\end{minted}
\caption{基本型の値を評価する}
\label{lst:ch01-basic-values}
\end{listing}

\begin{leanbox}
\texttt{List}と\texttt{Option}は、それだけでは具体的な値の型にならない。
\texttt{List Nat}、\texttt{List String}、\texttt{Option Nat}、\texttt{Option String}のように、内側に入る型を指定して使う。
このように、型から新しい型を作るものを、本章では型を受け取る構成子として読む。
\end{leanbox}

\section{\texttt{def}による名前付き定義}

Leanでは、\texttt{def}を使って名前付きの定義を書く。
定義できるものは、値でも関数でもよい。
まずは、構文の形を正確に読むことから始める。

\begin{center}
\texttt{def 名前 (引数 : 引数の型) : 返り値の型 := 本体}
\end{center}

\texttt{:=}の左側には、定義する名前、引数、返り値の型を書く。
\texttt{:=}の右側には、その定義の本体を書く。
Leanは、本体の型が返り値の型と合っているかを検査する。

\begin{listing}[htbp]
\begin{minted}{lean4}
def answer : Nat :=
  42

def addOne (n : Nat) : Nat :=
  n + 1

def add (m : Nat) (n : Nat) : Nat :=
  m + n

def joinWithSpace (left : String) (right : String) : String :=
  left ++ " " ++ right

#check answer
#check addOne
#check add

#eval addOne 5
#eval add 2 3
#eval joinWithSpace "Lean" "4"
\end{minted}
\caption{\texttt{def}による値と関数の定義}
\label{lst:ch01-def-basics}
\end{listing}

\texttt{answer}は引数を持たない定義である。
\texttt{answer : Nat}と書かれているので、\texttt{answer}は\texttt{Nat}型の値である。

\texttt{addOne}は一つの引数を持つ関数である。
\texttt{(n : Nat)}は、引数\texttt{n}が\texttt{Nat}型であることを表す。
返り値の型も\texttt{Nat}なので、\texttt{addOne}は\texttt{Nat}を受け取り\texttt{Nat}を返す関数である。

\texttt{add}と\texttt{joinWithSpace}は、二つの引数を持つ関数である。
Leanでは、複数の引数を\texttt{(m : Nat) (n : Nat)}のように並べて書く。
関数を呼び出すときも、\texttt{add 2 3}のように、関数名の後ろに引数を空白で並べる。

\begin{warningbox}
Leanの\texttt{def}は、可変変数への代入ではない。
\texttt{def answer : Nat := 42}は、「\texttt{answer}という名前を\texttt{42}という定義に結びつける」という意味である。
あとから同じ名前に別の値を代入して更新する、という読み方はしない。
\end{warningbox}

\section{\texttt{\#eval}による評価}

\texttt{\#eval}は、Leanに式を評価させるためのコマンドである。
\texttt{def}で定義した名前は、\texttt{\#eval}の中で使える。
Leanは定義を展開し、計算できる部分を計算し、結果を表示する。

\begin{listing}[htbp]
\begin{minted}{lean4}
#eval answer
#eval addOne answer
#eval addOne (addOne 10)
#eval add (addOne 2) (addOne 3)
#eval joinWithSpace (joinWithSpace "Lean" "4") "book"
\end{minted}
\caption{定義を使った評価}
\label{lst:ch01-eval-basics}
\end{listing}

\texttt{addOne (addOne 10)}では、内側の\texttt{addOne 10}が先に使われ、その結果に外側の\texttt{addOne}が適用される。
括弧は、どの部分を先にひとまとまりとして読むかを示している。

関数適用は、Leanでは空白で書く。
\texttt{addOne 5}は、\texttt{addOne}という関数を\texttt{5}に適用する式である。
二引数の関数も同じで、\texttt{add 2 3}と書く。

\begin{leanbox}
\texttt{\#eval}は、証明ではない。
一つの式を評価して、その結果を表示するコマンドである。
後の章で扱う\texttt{theorem}や\texttt{example}は、「すべての場合に成り立つ主張」をLeanに確認させるための構文であり、\texttt{\#eval}とは役割が異なる。
\end{leanbox}

\begin{warningbox}
\texttt{\#eval}は、Leanが計算でき、結果を表示できる式に対して使う。
型を確認したいときは\texttt{\#check}、式を評価したいときは\texttt{\#eval}、という使い分けを覚えるとよい。
\end{warningbox}

\section{関数型・関数適用・無名関数}

関数にも型がある。
\texttt{A -> B}は、\texttt{A}型の値を受け取り、\texttt{B}型の値を返す関数の型である。
たとえば、\texttt{Nat -> Nat}は自然数から自然数への関数の型である。

\texttt{def addOne (n : Nat) : Nat := n + 1}と書くと、\texttt{addOne}の型は\texttt{Nat -> Nat}になる。
\texttt{\#check addOne}で確認できる。

関数を名前なしで直接書くこともできる。
そのための構文が\texttt{fun}である。

\begin{center}
\texttt{fun 引数 => 本体}
\end{center}

\texttt{fun (n : Nat) => n + 1}は、「\texttt{n}を受け取り、\texttt{n + 1}を返す関数」をその場で作る式である。
このような関数を無名関数、またはラムダ式と呼ぶ。

\begin{listing}[htbp]
\begin{minted}{lean4}
#check addOne
#check (fun (n : Nat) => n + 1)

def inc : Nat -> Nat :=
  fun n => n + 1

def applyTwice (f : Nat -> Nat) (x : Nat) : Nat :=
  f (f x)

#eval inc 4
#eval applyTwice inc 10
#eval applyTwice (fun n => n + 2) 10
\end{minted}
\caption{関数型と無名関数}
\label{lst:ch01-function-types-lambda}
\end{listing}

\texttt{inc}は、返り値の型として\texttt{Nat -> Nat}を明示している。
右辺の\texttt{fun n => n + 1}では、\texttt{n}の型を書いていない。
これは、左辺に\texttt{Nat -> Nat}と書かれているため、Leanが\texttt{n : Nat}と推論できるからである。

\texttt{applyTwice}は、関数を引数として受け取る関数である。
第一引数\texttt{f}の型は\texttt{Nat -> Nat}である。
本体\texttt{f (f x)}は、まず\texttt{x}に\texttt{f}を適用し、その結果にもう一度\texttt{f}を適用する。

\begin{leanbox}
\texttt{A -> B -> C}という型は、Leanでは右に結合して\texttt{A -> (B -> C)}と読まれる。
また、\texttt{f x y}という関数適用は、左に結合して\texttt{(f x) y}と読まれる。
この二つの規則により、複数引数の関数を自然に書ける。
\end{leanbox}

\section{関数合成と型パラメータ}

二つの関数は、出力の型と入力の型が合うときにつなげられる。
\texttt{f : A -> B}と\texttt{g : B -> C}があるなら、まず\texttt{f}を適用し、その結果に\texttt{g}を適用することで、\texttt{A -> C}型の新しい関数を作れる。
これを関数合成と呼ぶ。

Leanでは、関数合成を自分で次のように定義できる。

\begin{listing}[htbp]
\begin{minted}{lean4}
def compose {A B C : Type} (g : B -> C) (f : A -> B) : A -> C :=
  fun x => g (f x)

def surround (s : String) : String :=
  "[" ++ s ++ "]"

def mark (s : String) : String :=
  s ++ "!"

def transform : String -> String :=
  compose mark surround

#eval transform "x"
#eval compose addOne addOne 10
\end{minted}
\caption{関数合成を定義する}
\label{lst:ch01-compose-definition}
\end{listing}

\texttt{compose}の定義を順に読む。
\texttt{\{A B C : Type\}}は、\texttt{A}、\texttt{B}、\texttt{C}が型であることを表す。
波括弧\texttt{\{...\}}で囲まれているので、通常はLeanがこれらを推論する。
たとえば\texttt{compose mark surround}では、\texttt{surround : String -> String}と\texttt{mark : String -> String}から、\texttt{A}、\texttt{B}、\texttt{C}はいずれも\texttt{String}だと分かる。

\texttt{(g : B -> C)}は、\texttt{B}から\texttt{C}への関数を受け取るという意味である。
\texttt{(f : A -> B)}は、\texttt{A}から\texttt{B}への関数を受け取るという意味である。
返り値の型は\texttt{A -> C}である。
つまり、\texttt{compose}は二つの関数を受け取り、新しい関数を返す。

本体\texttt{fun x => g (f x)}では、入力\texttt{x}にまず\texttt{f}を適用し、次に\texttt{g}を適用している。
\texttt{compose g f}という順番は、「先に\texttt{f}、後で\texttt{g}」と読む。
これは数学でよく使う合成記法の順番に近い。

\begin{intuitionbox}
関数合成で最も大切なのは、型がつながることを見ることである。
\texttt{A -> B}の出力\texttt{B}が、\texttt{B -> C}の入力\texttt{B}と一致しているから、二つの関数をつなげられる。
この見方が、第5章で扱う「射の合成」につながる。
\end{intuitionbox}

\FigurePlaceholder{fig-ch01-pipeline.png}{0.90}{関数合成による文字列変換}{fig:ch01-pipeline}

\section{サンプル：小さな文字列変換パイプライン}

最後に、ここまでに出てきた構文を一つの小さな例にまとめる。
新しい考え方を増やすのではなく、型注釈、\texttt{def}、\texttt{\#eval}、\texttt{Option}、\texttt{match}、\texttt{List.map}が同じコードの中でどう並ぶかを確認する。

\begin{listing}[htbp]
\begin{minted}{lean4}
def nonempty (s : String) : Option String :=
  if s == "" then none else some s

def decorate (s : String) : String :=
  mark (surround s)

def decorateIfNonempty (s : String) : Option String :=
  match nonempty s with
  | none => none
  | some value => some (decorate value)

def decorateAll (xs : List String) : List (Option String) :=
  xs.map decorateIfNonempty

#eval decorateIfNonempty ""
#eval decorateIfNonempty "lean"
#eval decorateAll ["a", "", "b"]
\end{minted}
\caption{文字列変換パイプライン}
\label{lst:ch01-string-pipeline}
\end{listing}

\texttt{nonempty}は、空文字列の場合に\texttt{none}を返し、それ以外の場合に\texttt{some s}を返す。
返り値の型が\texttt{Option String}なので、結果は「文字列がある場合」と「ない場合」に分かれる。

\texttt{decorateIfNonempty}では、\texttt{match}を使って\texttt{Option String}を場合分けしている。
\texttt{none}の場合は\texttt{none}を返す。
\texttt{some value}の場合は、取り出した\texttt{value}に\texttt{decorate}を適用し、もう一度\texttt{some}で包む。

\texttt{decorateAll}では、\texttt{xs.map decorateIfNonempty}と書いている。
ここでの\texttt{map}は、リストの各要素に同じ関数を適用し、結果をリストとして返す操作である。
\texttt{List.map}の性質は第8章で詳しく扱う。
本章では、ドット記法\texttt{xs.map f}を「\texttt{xs}の各要素に\texttt{f}を適用する」と読めれば十分である。

\section{本章でまだ扱わないこと}

本章では、Leanを「仕様を書くための言語」として使い始めた。
ただし、まだ扱わないことも多い。

\begin{warningbox}
本章では、\texttt{Prop}、\texttt{theorem}、\texttt{example}、\texttt{rfl}、\texttt{simp}、\texttt{rw}を本格的には扱わない。
これらは次章以降で導入する。
本章の目的は、証明を急ぐことではなく、Leanの型、値、関数、実行確認を読めるようになることである。
\end{warningbox}

また、本章の文字列変換の例は、言語構文を確認するための小さな例である。
文字列処理そのものの詳細には踏み込まない。
ここでは、Leanで関数の形、場合分け、リストへの適用を確認するために、意図的に単純なモデルにしている。

\section{実務への読み替え}

本章の内容は、すぐに次のような実務上の読み替えができる。

\begin{itemize}
\item 入力検査は、\texttt{String -> Option String}のように失敗可能性を型へ出せる。
\item DTO変換は、\texttt{A -> B}という関数として小さく切り出せる。
\item 複数件の変換は、\texttt{List A -> List B}や\texttt{List A -> List (Option B)}として表せる。
\item パイプラインは、関数合成として分解して読める。
\end{itemize}

ここで大切なのは、Leanで実務システム全体を再実装することではない。
実務コードの中にある「仕様上重要な変換」を小さな関数として切り出すことである。
小さく切り出せると、次章以降で「その関数はどんな性質を満たすべきか」を命題として書けるようになる。

\section{圏論への橋渡し}

本章では、圏論という言葉をほとんど使わなかった。
しかし、すでに圏論へ進むための材料は出そろっている。

\begin{center}
\begin{tabular}{lll}
\toprule
本章で見たもの & 後の圏論での読み方 & ソフトウェアでの感覚 \\
\midrule
型 & 対象 & データの形、状態空間 \\
関数 & 射 & 変換、処理、adapter \\
関数合成 & 射の合成 & パイプライン、処理連結 \\
何もしない関数 & 恒等射 & no-op、恒等変換 \\
\texttt{List.map} & 関手の入口 & 構造を保った一括変換 \\
\texttt{Option} & 文脈つきの値 & 失敗可能性、nullable \\
\bottomrule
\end{tabular}
\end{center}

第5章では、これらを「合成できる処理の世界」として整理する。
その前に、第2章から第4章で、命題、証明、等式、不変条件を扱う。
型と関数だけでは、「この変換が正しい」とまでは言えない。
正しさを言うには、関数が満たすべき性質を命題として書く必要がある。

\section{本章のまとめ}

Leanでは、型を値に対する約束として読むとよい。

\texttt{Nat}、\texttt{Bool}、\texttt{String}、\texttt{List}、\texttt{Option}は、仕様の小さな部品を表すための入口になる。

\texttt{def}は、名前付きの値や関数を定義するための基本構文である。

\texttt{\#check}は型を確認し、\texttt{\#eval}は計算できる式を評価する。

\texttt{Option}を使うと、値が存在しない可能性を型に明示できる。

関数合成は、処理パイプラインとして読める。

本章の内容は、後の命題・証明・圏論の基礎になる。
